\documentclass[10pt]{article}
\usepackage[margin=20mm]{geometry}
\usepackage{amsmath,amssymb,fontspec,unicode-math,xurl,hyperref}
\IfFontExistsTF{DejaVu Serif}{\setmainfont{DejaVu Serif}}{\setmainfont{Latin Modern Roman}}
\setmathfont{latinmodern-math.otf}
\setlength{\parindent}{0pt}\setlength{\parskip}{6pt}\sloppy
\begin{document}
Date:\par
\# Scientific verification report: I25-40165 (0058)\par
Supported impact: Major, scoped to two central quantitative derivations.
Batch B001 / [internal task]. Token: [internal identifier omitted]
Scientific decision: pending human review. Author adoption has not been requested.\par
\#\# 1. Identity, version and inspected scope\par
Yash Goel, Ayan Sengupta and Tanmoy Chakraborty, Position: Enough of Scaling LLMs! Lets Focus on Downscaling, ICML 2025, PMLR 267:81375-81392 [E1-E2]. The primary reviewed object is the 18-page publisher-linked PDF. All page references are one-based physical PDF pages, matching the printed article pages.\par
Primary PDF SHA-256:
`f7645a566ade49e81a277e7e47ad9a69a9ebbacd136c2f8f0c6e56ef3d5fe9f0`\par
Its bytes match the controller acquisition receipt and remain preserved. The review read the main argument (pp. 1-9), all Appendices A-E (pp. 16-18), and selected cited primary sources. Source pages 4, 7, 16, 17 and 18 were also inspected visually. The publisher-linked repository was acquired at commit `c5c328c359456abbd512bc5d99885315b9a851fb`, dated 2 May 2025, 03:58:47 UTC. Its 11-file tree is preserved in `sources/0058/repository/` [E3].\par
The arXiv record lists v1 on 2 May, v2 on 5 May, and v3 on 25 May 2025. A separate v3 PDF was acquired; targeted comparison confirms the same factor deletion in Appendix E and the same carbon-performance equations and percentage claim. This is not a complete version diff, and the publisher PDF remains the reviewed authority [E4].\par
\#\# 2. Summary and scientific dependence\par
The paper advances a research position favoring smaller, specialized models. It also presents an explicit sufficient condition for an ensemble of pruned models to improve loss at unchanged modeled compute, and a quantitative relationship between downstream performance and carbon emissions. These technical claims supply substantive support for the position.\par
Two derivations fail even when their surrounding fitted-law assumptions are granted. First, Proposition 4.1 deletes a multiplicative pruning-loss factor that does not cancel. A complete finite countermodel satisfies the printed sufficient condition and the stated compute equality, yet has higher ensemble loss. Section 4.5 explicitly relies on this proposition to prescribe the proposed pipeline. Second, Eqs. (8)-(10) discard an affine intercept and reverse an exponent sign. The stated universal carbon cost of a 10\% performance increase does not follow.\par
These are Major findings because they damage central mathematical guarantees or quantitative conclusions. Neither finding requires proving that every small-model ensemble performs poorly, that the environmental motivation is false, or that the whole downscaling agenda should be rejected. The conditional fixed-hardware carbon proportionality and the elementary ensemble reparameterization have useful valid content, discussed below.\par
\#\# 3. Definitions and operative assumptions\par
Let N0 be the pre-pruning model size, D the post-training token count, rho the removed fraction, L0 the original loss, and n the number of equal-size pruned models. The derivation uses rho = (n-1)/n, so n must be at least 2 to avoid the singular no-pruning endpoint. Assume positive N0, D and b, and a common loss scale.\par
Define the factor retained in Eq. (11):\par
`Q = NC/N0\textasciicircum{}alpha + DC/D\textasciicircum{}beta + E.`\par
The two fitted relations used by Appendix E are:\par
- Pruned-model loss: `Lc = L0 + Q/(rho\textasciicircum{}gamma N0\textasciicircum{}delta).`\par
- Ensemble loss: `Le = Lc - b + b/n\textasciicircum{}a.`\par
The review grants these relations exactly for the algebraic countermodel. It does not claim that arbitrary coefficients describe a measured LLM. The paper additionally assumes Q is approximately constant at matched budget; the countermodel makes Q exactly constant. Equal modeled compute follows from `n K (1-rho) N0 D = K N0 D`, with fixed positive K. This does not by itself establish equal latency, energy, or orchestration overhead.\par
For the second finding, write carbon emissions as Ccarb to avoid confusing them with the paper’s computational budget C. The review grants its fixed-data, fixed-hardware simplification `Ccarb = Kc L\textasciicircum{}(-1/alpha)`, with Kc positive and alpha = 0.08, and its affine performance model `P = w1 + w2 L`.\par
\#\# 4. Problem 1: the sufficient downscaling condition drops a decisive factor\par
Judgment: Major, affecting Proposition 4.1 and the pipeline’s claimed guarantee.\par
\#\#\# Source claim and invalid step\par
Section 4.4, p. 7, presents Eq. (13) as a sufficient condition for strictly lower ensemble loss at the same compute. Appendix E, pp. 17-18, begins with the complete pruning law in Eq. (27), then replaces Q by 1 in Eq. (28). The explanation says Q is approximately constant and therefore cancels from both sides of the comparison. Section 4.5, p. 8, expressly directs pruning to obey this condition and uses it to assert better performance at equal cost [E1].\par
Substituting the two full fitted laws actually gives:\par
`Le - L0 = Q/(rho\textasciicircum{}gamma N0\textasciicircum{}delta) - b(1 - n\textasciicircum{}(-a)).`\par
Only L0 cancels. Q multiplies the pruning penalty and has no identical additive counterpart on the right-hand side. Its being constant does not make its magnitude irrelevant.\par
Let `g(n) = (1 - n\textasciicircum{}(-a)) ((n-1)/n)\textasciicircum{}gamma`. The printed condition is `g(n) >= 1/(b N0\textasciicircum{}delta)`. Under the assumed laws, the correct non-worse condition is:\par
`g(n) >= Q/(b N0\textasciicircum{}delta).`\par
Strict improvement requires a strict inequality. The printed non-strict antecedent cannot guarantee strict improvement even in the special case Q = 1.\par
\#\#\# Complete countermodel\par
Choose the following exact values:\par
Quantity | Value\par
n, N0, D, L0 | 2, 8, 1, 2\par
a, b, gamma, delta, alpha, beta | all 1\par
NC, DC, E | 8, 2, 2\par
rho | 1/2\par
Q | 8/8 + 2/1 + 2 = 5\par
Every size and token-count variable is positive; n is a positive integer greater than 1; the pruning fraction lies strictly between 0 and 1; all fitting coefficients here are positive; and Q is exactly fixed. The example is an algebraic countermodel to the claimed implication of the assumed laws, not a purported empirical eight-parameter language model.\par
The printed condition passes strictly:\par
`g(2) = (1 - 1/2)(1/2) = 1/4 > 1/8.`\par
But the unabridged relations give:\par
`Lc = 2 + 5/[(1/2) 8] = 13/4,`\par
`Le = 13/4 - 1 + 1/2 = 11/4 > L0 = 2.`\par
The modeled compute is unchanged: the two retained-size-4 models each use D = 1, giving total cost 8K, equal to the base model’s 8K. Thus neither a violated budget nor a boundary equality explains the failure. The correct condition fails, as it should: 1/4 is less than 5/8. Exact rational evaluation is preserved with the independent mathematical audit [E5].\par
\#\#\# Alternatives, counterevidence and repair\par
One could introduce an explicit normalization Q = 1, or redefine the ensemble coefficient after rescaling the loss. Neither follows from Q being approximately constant; either changes the hypotheses or coefficients. Similarly, allowing approximation error does not establish a robust sufficient condition without a bound on the omitted factor and on the two fitted laws.\par
The p. 7 LLaMA illustration also uses coefficients transferred across settings. Upstream Chen arXiv:2411.10272v2 prints gamma = +1.08 in the relevant fitted table, while v3 changes it to -1.08. The target’s positive sign can therefore be faithful to the earlier cited version and is not classified as an independent misquotation. The retained evidence preserves this history; the demonstrated factor omission is independent of it. No claim that the printed threshold is the smallest possible n is needed.\par
A precise conditional repair is to retain Q, use the corrected threshold above, and distinguish non-worse from strictly better loss. This repairs the algebra exactly. It does not validate transferring CNN ensemble coefficients to domain-specialized, pruned LLMs. Such transfer requires a consistent loss, units, pruning-law version, ensemble diversity assumptions and suitable measurements. If an 8B model is divided under this compute model, each member has size 8B/n; only n = 8 gives 1B members at the same total size.\par
\#\# 5. Problem 2: the carbon-performance conversion is algebraically invalid\par
Judgment: Major, affecting the claimed quantitative environmental-performance relationship.\par
\#\#\# Source claim and exact derivation\par
On p. 4, Eq. (7) relates carbon to loss with exponent -1/alpha. Eq. (8) states `P = w1 + w2 L`. Eq. (9) then uses carbon proportional to `P\textasciicircum{}(-1/alpha)`, and Eq. (10) changes to performance proportional to carbon raised to positive alpha. The text derives a fixed carbon multiplier for a 10\% performance gain [E1].\par
Even granting Eq. (7), correct substitution gives:\par
`Ccarb = Kc ((P - w1)/w2)\textasciicircum{}(-1/alpha),`\par
on the domain `(P - w1)/w2 > 0`, or equivalently:\par
`P = w1 + w2 (Ccarb/Kc)\textasciicircum{}(-alpha).`\par
The additive intercept cannot be absorbed into a proportionality constant. Moreover, even if w1 were zero, inverting Eq. (9) gives a negative exponent, not the positive exponent in Eq. (10). Redefining performance as inverse loss would be a different premise, inconsistent with the displayed affine model.\par
The cited Chen et al. primary source fits an affine loss-to-performance relation to emerged, above-random observations. It neither sets the intercept to zero nor establishes a universal positive-exponent carbon law. Its current v2 also explicitly discusses limits of composing compute and performance fits across different training trajectories [E6, p. 3, Sections 2.2-2.3].\par
\#\#\# Complete numerical check\par
Choose alpha = 0.08, w1 = 1 and w2 = -1/2. Let performance increase from 0.50 to 0.55, exactly 10\% relative. The corresponding positive losses are L1 = 1 and L2 = 0.9. Both performance values lie in [0,1], and the sign expresses the intended improvement as loss decreases.\par
These values can be realized within the paper’s assumed scalar laws by taking a positive reference size Nr, setting `L = (N/Nr)\textasciicircum{}(-0.08)`, holding D and hardware fixed, and taking carbon proportional to N. Thus the witness satisfies the premises used in the carbon substitution.\par
Calculation | Carbon ratio C2/C1\par
Correct consequence of Eqs. (7)-(8): 0.9\textasciicircum{}(-12.5) | 3.732232\par
Printed Eq. (9): 1.1\textasciicircum{}(-12.5) | 0.303803\par
Printed Eq. (10): 1.1\textasciicircum{}(12.5) | 3.291611\par
The three outcomes disagree. This is a check of the claimed implication, not a new measured carbon estimate. More generally, the ratio depends on the intercept and starting performance; there is no task-independent 10\% rule from the stated premises. The reviewer-authored scalar calculation and output are preserved [E7].\par
The displayed arithmetic also confuses percentage of baseline with percentage increase: 3.291611 times the baseline is 329.161\% of the baseline, or 229.161\% more. This local arithmetic issue is subordinate to the invalid model conversion. A fixed power law is also not an exponential law in performance.\par
\#\#\# Impact and precise repair\par
This removes the stated quantitative carbon-versus-performance conclusion used in the paper’s environmental argument. It does not establish that scaling is environmentally benign. Retaining the affine offset and exponent gives the valid conditional relation above; estimating a practical multiplier additionally requires task-specific coefficients, an operating point, the scaling trajectory for N and D, and appropriate hardware assumptions. No corrected universal percentage is justified.\par
\#\# 6. Other checked claims and evidence limits\par
Conditional carbon proportionality. Given the paper’s fixed hardware, power, efficiency, lifetime and data-center coefficients, substituting runtime proportional to ND into its emission formulas does yield carbon proportional to ND. This is a conditional algebraic result. Changes in those coefficients or in the scaling path require separate treatment. The proof does not make K1 + K2 universal across deployments.\par
What Figures 2-3 show. Static inspection of `analysis.ipynb` identifies a single GPT3 database row in zero-based cell 11. Cells 18-19 vary N or D while reusing its hardware settings, append predicted emissions, and cells 20-21 plot the resulting curves. `model.py`, lines 68-117, computes FLOPs from 6ND and carbon from power times modeled runtime. Cells 14-17 similarly evaluate analytic loss bounds. These are model-generated curves; the figure labels themselves say estimated. A cached single-row residual is not independent validation of the across-size scaling relationship. The p. 4 statement of empirical validation therefore remains unsupported by an independent scaling experiment in the inspected release [E3]. This qualification does not invalidate the conditional algebra.\par
Ensemble reparameterization. Appendix D correctly derives Eq. (12) from L = c + b/n\textasciicircum{}a by setting the single-member loss to c + b, provided the member family and fitted coefficients remain fixed. That valid rearrangement does not justify the subsequent deletion of Q or guarantee transfer to heterogeneous domain specialists.\par
Loss-floor interpretation. A finite-scale lower bound depending on N and D is not a positive bound uniform in both variables. In the displayed Kaplan formula with positive exponents, both terms tend to zero when N and D grow together. Chinchilla’s explicit positive E can supply a floor. The p. 4 statement should keep these cases distinct; the review does not infer that real-world cross-entropy can be driven to zero.\par
The recommendations for data curation, pruning and domain specialization were read as research proposals. Their merits are not rejected for lacking a new benchmark experiment. Conversely, the empirical applicability of the proposed composed law remains unverified. The review did not reproduce training, inference, carbon measurements, the complete cited literature or every historical source revision.\par
\#\# 7. Conclusion and remaining work\par
Item | Judgment | Repair status | Adoption / remaining work\par
Proposition 4.1 drops Q | Major | Correct conditional threshold proved | Not applied; refit and validate transfer\par
Carbon-performance conversion | Major | Affine relation restored analytically | No universal multiplier; coefficients needed\par
Figures 2-3 empirical validation | Needs Evidence | Model-based scope established statically | Independent measurements absent from inspected release\par
Fixed-coefficient carbon algebra; Appendix D | No Issue Found in checked steps | Not applicable | Keep stated conditions\par
Broad downscaling recommendation | Not refuted by these checks | Not a mathematical repair | Human scientific assessment pending\par
The original publisher PDF and the separately saved arXiv v3 preserve the inspected claims. No full errata search or author correspondence was performed. Future experiments are necessary for the proposed LLM guarantee’s applicability; they are not necessary to establish the invalid algebra already demonstrated.\par
\#\# 8. References and evidence\par
- E1: Exact target publisher PDF, `papers/0058/0058.pdf`; extraction `sources/0058/paper.layout.txt`; identity receipt and source-page renders.\par
- E2: https://proceedings.mlr.press/v267/goel25c.html ; saved publisher snapshot.\par
- E3: https://github.com/LCS2-IIITD/Downscaling/tree/c5c328c359456abbd512bc5d99885315b9a851fb ; preserved repository and `verification/0058/code\_audit/` acquisition manifest, numbered excerpts and audit notes.\par
- E4: https://arxiv.org/abs/2505.00985 ; v3 dated 25 May 2025; separate PDF SHA-256 `9b2c5d3f6787f4fe7a9447c379dc4f36ba18f0ac27738829cb13688703f88994`. Targeted comparison saved in `arxiv\_v3\_targeted\_comparison.json`.\par
- E5: `verification/0058/math\_audit/`: independent calculations, audit notes and versioned primary-source evidence, including Chen et al., P2 Law, arXiv:2411.10272; Lobacheva et al., On Power Laws in Deep Ensembles, arXiv:2007.08483v2.\par
- E6: Chen et al., Scaling Laws for Predicting Downstream Performance in LLMs, arXiv:2410.08527v2, 7 April 2025, accepted to TMLR. https://arxiv.org/abs/2410.08527 ; saved PDF SHA-256 `8cc5914aedbb7b4175eba5a9832aa62aacafe2fd8460f7e2749105785c8e416b`.\par
- E7: `verification/0058/root\_independent\_checks.py` and `root\_independent\_checks.json`.\par
- Workflow: Paper 01 / Paper 02 evidence guidance. Kassis et al. (2026), Scientific Agent Skills, https://doi.org/10.48550/arXiv.2609.00065 ; procedural reference, not evidence for the scientific findings.\par
\#\# Appendix: execution, provenance and roles\par
The repository archive is 112,490 bytes, SHA-256 `4035fb164879ab5bd86b52dea6f70f00bf0ba3f6ea4539dfa5bbfa6e9eec911f`. Author files and notebook outputs were treated as data. No author module, notebook or training workflow was executed or imported. No software was installed, no paid model/API was used, and no author source or controller file was changed.\par
Existing Poppler extracted and rendered sources. Reviewer-authored scripts performed small scalar and rational calculations and wrote evidence records. Existing Pandoc converted this report to standalone TeX. The initial PDF used the existing Markdown renderer. The accept-0012 replacement PDF is compiled from this TeX with the existing offline Tectonic runtime. Exact compilation and page-inspection hashes are recorded in the acceptance files.\par
The primary assistant conducted the review with independent AI passes on mathematical reasoning and static code evidence. No human review or author approval is claimed. Provenance and coverage records distinguish source reading, static inference, independent calculation, and work not performed. Scientific decisions and controller acceptance remain pending.\par
\end{document}
